% Transcription — fonds Alexandre Grothendieck, shelfmark 14, batch 2 (pages 21-40).
% Established from the facsimile archives/batches/14/batch-02.pdf (a mirror of
% https://grothendieck.umontpellier.fr/14.pdf), per the skill
% transcribe-grothendieck. The batch file carries no cover sheet: PDF page N
% is page 20+N of the folder (the Montpellier footer prints « 21 » ... « 40 »),
% and the archivists' pencil numbers bottom-left agree leaf by leaf (21 to 38
% and 40 are legible; 39 is faint but in sequence).
%
% Pass: Opus 5.5 (claude-opus-5-5), 2026-09-23 — first pass, unchecked against the pages by a human.
%
% Seventeen of the twenty pages are transcribed. Absent: page 21, a leaf of
% an unrelated typescript (a Bourbaki draft, « n° 391 », p. 14, on affine
% pseudo-reflections and Proposition 9 on orthogonal symmetries, scanned
% upside down, nothing in his hand); page 22, blank; page 39, an otherwise
% blank sheet inscribed in his hand « Astuce de Tate sur les cycles
% décomposables » — a cover, whose words become the \section heading before
% page 40, per references/hand.md §5.
%
% What the batch is. Two runs in his hand.
%   (1) Pages 23-38, headed « Théorie de l'équivalence d'Albanese des
%   cycles », written fairly with numbered sections, definitions and
%   theorems, not paginated by him:
%   23     §1 « L'anneau des classes d'équivalence d'Albanese ». Def. 1.1:
%          a cycle z on X smooth over k is alb-equivalent to 0 if z = Z(t)
%          for a cycle Z on X×T, T smooth connected, t a 0-cycle of degree 0
%          on T with S_T(t) = 0 (Albanese sum). Prop. 1.2 (i): the subgroup
%          N(X), containing linear equivalence, image N'^*(X) in the Chow
%          ring A^*(X).
%   24     (ii) N'^* is an ideal, ring A'^*(X); (iii) f^*; (iv) f_* (Gysin);
%          (v) Albanese equivalence implies algebraic equivalence. Proof of
%          (i) begins: z - z' via T×T'.
%   25-27  the proof of (i) (margin « Lemme à dégager »), (iii) via
%          (f×id_T)^*; (ii) by the same method, (iv) by « l'AQT » (?), (v)
%          trivial. Scholie 1.3: A'^* has the formal properties of A^*;
%          equal to A^* where linear = algebraic equivalence (P^n, flag
%          varieties).
%   28-29  Th. 1.4: for 0-cycles, alb-equiv. to 0 iff degree 0 and S_X = 0;
%          for divisors, iff linearly equivalent to 0 (via u: A_T → Pic^0_X).
%          Rem. 1.5: unknown, already for 0-cycles, whether Albanese
%          equivalence is linear equivalence. Th. 1.6: products of classes
%          algebraically equivalent to 0 vanish in A'^*(X); (x)-(y) × (x')-(y')
%          has S_{T×T'} = 0 (margin: not known to be lin. equiv. to 0, even
%          for curves).
%   30-33  §2 « Familles de classes de cycles mod équiv. d'Alb. ». Th. 2.1:
%          Z ∈ A^i_alb(X×Y) gives u_Z : A_Y(k) → A^i_alb(X), (i) depending only
%          on the τ-class of Z, (ii) up to isogeny A' → A_Y induced by a
%          homomorphic family Z'. Proof via a cycle U on Y×A_Y with
%          S_Y(U(a)) = na and divisibility of A_Y(k). Cor. 2.2 (algebraic
%          equivalence = Albanese-equivalent to a family from an abelian
%          variety), 2.3 (P^i(X) is divisible), 2.4 (alg. × τ-equivalent to
%          0 has product 0).
%   34-36  Def. 2.5 algebraic maps Y(k) → A^i_alb(X); Prop. 2.6, 2.7
%          (functoriality). N.B.: reduction to algebraic homomorphisms
%          A(k) → P^i(X), P^j(X') → P^i(X); duality between P^i(X) and P^j(X),
%          i+j = n+1, for dual abelian varieties A, B; question of transposition.
%   37-38  §3 « Homomorphismes ℓ-adiques associés à une famille abélienne de
%          cycles »: Th. 3.1 u_Z = 0 ⟹ a(tZ) = 0 on H^{2(n-i)+1}(X)(n-i) → H^1(A)
%          (« dém. pas faite »; « A-t-on une réciproque ? »); reduction to an
%          ample curve C, Künneth decomposition, the middle factor. Page 38
%          stops two-thirds down; the run is not concluded here.
%   (2) Page 40, after the cover « Astuce de Tate sur les cycles
%   décomposables », in a darker ink on yellowed paper: A an abelian variety
%   over k algebraically closed, T_ℓ(A) a module over the « pseudo-algèbre de
%   Lie » g; hypotheses a) Tate's conjecture for H^2(A, Q_ℓ(1)) and b) g_K
%   « pseudo-algébrique » or T_K(A) = V_1 × V_2; verified for a power of an
%   elliptic curve over F_q, a) « vérifié par Mumford ». The page ends
%   mid-sentence, « (comme on sait », continued presumably in batch 3.
%
% Notation fixed for the batch. \mathfrak{z}, \mathfrak{z}' for his curly z
% (a cycle on X, the ʒ of hand.md §4); Z, Z' for the cycles on X×T;
% \mathfrak{t} for the 0-cycle on T (a figure like « Vl », read as a fraktur
% t, for T); S_T, S_X, S_Y the Albanese sum; A^*(X) the Chow ring, A'^*(X) its
% quotient, A^i_{\mathrm{alb}}(X) (he writes « alb » as a subscript from
% page 30 on), A_Y the Albanese variety of Y, A^*_Y its « torseur » of all
% degrees as he writes it; P^i(X) the classes algebraically equivalent to 0;
% \mathbb{T}_\ell(A), \mathfrak{g} on page 40. His « alb-équiv. », « lin.
% équiv. », « alg. équiv. », « s-groupe », « hom », « VA », « ssi » are kept.
% Plain square brackets in the text are his: either his own brackets, or
% words he added in the interline (« [noté N(X)] », « [connexe] »); the
% editorial additions are \add{} and none was needed.
%
% Legibility. The statements carry; the fast connective prose of pages 30-32,
% 36 and 40 and the rotated margins of pages 28 and 29 are partly \ill{}.
% Nothing on these leaves is dated, and the letter (1965) of the folder is not
% in this batch.
%
% The dating is the inventory's own for the folder, copied verbatim. Its
% group, [10-18] « Théorie des motifs », is dated [à partir de 1958-à partir
% de 1982].
\documentclass[11pt,a4paper]{article}
% Le préambule commun à toutes les transcriptions du fonds Grothendieck.
%
% Il tient en une idée : **l'incertitude se compose, elle ne se résout pas.**
% Une transcription qui lisse un mot douteux a détruit la seule chose pour
% laquelle on la faisait — un lecteur doit pouvoir savoir, à chaque endroit,
% ce qui était sur le feuillet et ce qui a été ajouté. Les six macros qui
% suivent sont donc l'appareil critique, et non de la décoration.
%
% Les mêmes commandes sont reconnues par `scripts/render.mjs`, qui produit la
% vue de lecture du site. Une macro ajoutée ici doit l'être aussi là-bas,
% faute de quoi le rendu échouera bruyamment — ce qui est voulu : un rendu
% silencieusement faux est pire qu'un rendu absent.

\ProvidesPackage{grothendieck}[2026/08/08 Transcriptions du fonds Grothendieck]

\usepackage{amsmath,amssymb,amsthm}
\usepackage{mathrsfs}
\usepackage{stmaryrd}
% Les diagrammes commutatifs sont partout dans ce fonds : c'est la langue
% même de la géométrie algébrique de Grothendieck, et une transcription qui
% les rendrait en prose ne transcrirait rien.
\usepackage{tikz-cd}
% Français par défaut — c'est la langue des feuillets — mais l'anglais est
% chargé aussi : la traduction et le résumé sont des documents anglais, et la
% typographie française (espaces avant les deux-points) y serait une faute.
% Ces documents basculent par \selectlanguage{english} en tête de corps.
\usepackage[english,french]{babel}
\usepackage{fontspec}
\usepackage{geometry}
\geometry{a4paper,margin=25mm,marginparwidth=28mm}
\usepackage{marginnote}
\usepackage{xcolor}
% ulem, not soul. `soul`'s \st and \ul are built on a character-by-character
% scan that XeLaTeX's font machinery breaks: the document fails with an
% undefined control sequence rather than degrading. ulem does the same two jobs
% and is Unicode-engine safe. `normalem` stops it hijacking \emph, which the
% transcriptions use for Grothendieck's own emphasis.
\usepackage[normalem]{ulem}
\usepackage{hyperref}
\hypersetup{colorlinks=true,linkcolor=black,urlcolor=black,pdfborder={0 0 0}}

% --- Métadonnées du lot -----------------------------------------------------
% Reprises telles quelles par le script de rendu, qui les affiche en tête de
% la vue de lecture. Elles fixent ce que le fichier prétend couvrir : sans
% elles, une transcription flotte sans point d'attache dans un fonds de seize
% mille feuillets.

\newcommand{\folder}[1]{\gdef\grothFolder{#1}}
\newcommand{\batch}[1]{\gdef\grothBatch{#1}}
\newcommand{\leaves}[2]{\gdef\grothFirst{#1}\gdef\grothLast{#2}}
\newcommand{\foldertitle}[1]{\gdef\grothFoldertitle{#1}}
\gdef\grothFolder{?}\gdef\grothBatch{?}\gdef\grothFirst{?}\gdef\grothLast{?}\gdef\grothFoldertitle{}

% --- Le résumé --------------------------------------------------------------
%
% Il ouvre la lecture modernisée, sous le titre « Résumé », et il est composé
% en petit. Ce n'est pas un ornement : c'est le seul passage du document qui
% ne soit pas des mathématiques, et un lecteur doit pouvoir le reconnaître
% avant de l'avoir lu — puis le sauter s'il connaît déjà le sujet, ou s'y
% arrêter s'il ne le connaît pas. Le filet qui le suit marque la reprise du
% corps.

\newenvironment{resume}
  {\small\setlength{\parskip}{0.45em}}
  {\par\medskip\hrule\medskip}

% --- Mots-clés ---------------------------------------------------------------
%
% En anglais, en clôture du résumé de la lecture modernisée : le vocabulaire
% moderne sous lequel chercher ce que les feuillets construisent. Le manifeste
% du site (`npm run manifest`) les extrait pour étiqueter la cote entière —
% cette ligne est la source unique des tags, il n'y a pas de fichier à côté.
\newcommand{\keywords}[1]{\par\smallskip\noindent
  {\footnotesize\textsc{Keywords} --- \textit{#1}}\par}

% --- L'appareil critique ----------------------------------------------------

% Le numéro de feuillet, en marge. C'est l'ancre qui permet de régler un
% désaccord sur une formule en regardant *un* feuillet plutôt qu'une chemise
% de six cents. Le site s'en sert aussi pour tourner les pages du fac-similé
% quand on fait défiler la transcription.
\newcommand{\leaf}[1]{%
  \marginnote{\footnotesize\color{gray!70}\textsf{#1}}%
  \label{leaf:#1}\ignorespaces}

% Illisible : on ne devine pas. Un mot inventé qui se lit comme les autres est
% le pire résultat possible.
\newcommand{\ill}{\textcolor{red!60!black}{[\,\ldots\,]}}

% Lecture douteuse : le mot est donné, et signalé comme incertain.
\newcommand{\uncertain}[1]{\uline{#1}}

% Ajout de l'éditeur — une lettre manquante, un mot sous-entendu.
\newcommand{\add}[1]{\textcolor{blue!50!black}{[#1]}}

% Biffé par Grothendieck lui-même. Ce qu'il a rayé fait partie du document :
% une rature dit souvent où la pensée a changé de direction.
\newcommand{\struck}[1]{\sout{#1}}

% Note du transcripteur, sur l'état matériel du feuillet.
\newcommand{\note}[1]{\par\smallskip{\small\color{gray!60!black}\textit{[#1]}}\par\smallskip}

% Note marginale de Grothendieck, distincte du corps du texte.
\newcommand{\marginal}[1]{\par\smallskip\noindent\hspace{1em}%
  {\small\color{gray!60!black}#1}\par\smallskip}

% --- Titre ------------------------------------------------------------------

\AtBeginDocument{%
  \begin{center}
    {\large\bfseries Fonds Alexandre Grothendieck --- cote n\textsuperscript{o}~\grothFolder}\\[2pt]
    {\itshape\grothFoldertitle}\\[4pt]
    {\small feuillets \grothFirst--\grothLast\ (lot \grothBatch)}\\[2pt]
    {\footnotesize\color{gray!60!black}Transcription établie d'après le fac-similé numérisé
     par l'Université de Montpellier.\\
     Les lectures incertaines sont soulignées, les illisibles notées [\,\ldots\,],
     les ajouts entre crochets.}
  \end{center}
  \vspace{1em}}

\endinput


\folder{14}
\batch{2}
\pages{21}{40}
\dating{1965-[vers 1971]}
\watermark{Édition de démonstration}
\foldertitle{Théorie des cycles algébriques : conjectures de Hodge, Tate, Lefschetz, Weil : notes manuscrites (s.d.), lettre (1965).}

\begin{document}

\section{Théorie de l'équivalence d'Albanese des cycles}

\note{titre de sa main, souligné, en tête du feuillet 23. La rédaction qui
suit (feuillets 23 à 38) n'est pas paginée par lui ; elle est divisée en
paragraphes numérotés 1., 2., 3., soulignés, qui deviennent ici des
sous-sections}

\page{23}
\marginal{Consulter Weil, \uncertain{Critères} d'équivalence}

\subsection*{1. L'anneau des classes d'équivalence d'Albanese}

\note{« classes d' » est ajouté dans l'interligne, au-dessus d'un premier
« des » biffé}

\emph{Définition 1.1.} Soit $X$ un schéma \struck{projectif} lisse sur le
corps $k$, $\mathfrak{z}$ un cycle sur $X$. On dit que $\mathfrak{z}$ est
alb-équiv.\ à $0$, si on peut trouver un schéma [projectif] lisse
\struck{\ill{}} [connexe] $T$ sur $k$, \struck{un $0$-cycle de degré $0$}
$\mathfrak{t}$ sur $T$, et un cycle $Z$ sur $X \times T$, tels que
\marginal{?}
\begin{enumerate}
  \item[(i)] [$\mathfrak{t}$ est de degré $0$ et] $S_{T}(\mathfrak{t}) = 0$
  \item[(ii)] $Z(\mathfrak{t})$ est défini et égal à $\mathfrak{z}$.
\end{enumerate}
\note{les crochets autour de « projectif » sont de sa main ; « de degré 0
et » est ajouté au-dessus de (i). Le cycle sur $T$ est un signe semblable à
« Vl », transcrit $\mathfrak{t}$ dans tout le lot ; $S_T$ est la somme
d'Albanese}

\emph{Proposition 1.2.} (i) Les cycles alb-équiv.\ à $0$ sur $X$ forment un
sous-groupe [noté $N(X)$] (de $\mathrm{Cyc}^{\bullet}(X)$) contenant le
groupe des cycles lin.\ équiv.\ à $0$, de sorte que la connaissance de
$N^{*}(X)$ équivaut à celle de son image $N'^{*}(X)$ dans
\struck{$\mathrm{Cyc}$} $A^{*}(X)$ (l'anneau de Chow de $X$, défini par les
classes à équiv.\ linéaire près), et
\[
  \mathrm{Cyc}(X)/N(X) \simeq A^{*}(X)/N'^{*}(X).
\]

\page{24}
(ii) $N'^{*}(X)$ est un \emph{idéal} de $A^{*}(X)$, donc
$\mathrm{Cyc}^{*}(X)/N^{*}(X) \simeq A^{*}(X)/N'^{*}(X)$ est muni d'une
structure d'\emph{anneau}, noté $A'^{*}(X)$.

(iii) Pour tout hom $f : X \to Y$ de \struck{variétés projectives}
\uncertain{q.p.} lisses, l'hom
\[
  f^{*} : A^{*}(Y) \to A^{*}(X)
\]
applique $N'^{*}(Y)$ dans $N'^{*}(X)$, donc définit par passage au quotient
un hom d'anneaux $A'^{*}(Y) \to A'^{*}(X)$.

(iv) Pour tout $f$ comme dans (iii) \marginal{avec $f$ propre}, l'hom de
Gysin $f_{*} : A^{*}(X) \to A^{*}(Y)$ \struck{est} applique $N'^{*}(X)$ dans
$N'^{*}(Y)$, et induit donc un hom de groupes (également noté $f_{*}$)
\[
  A'^{*}(X) \to A'^{*}(Y)
\]

(v) L'équiv.\ d'Albanese $\Rightarrow$ équivalence algébrique. (Donc on peut
parler d'éléments de $A'^{*}(X)$ \emph{alg.}\ \emph{équiv.}\ à $0$ !)

\emph{Dém.} (i) \struck{Si $\mathfrak{z}$ \uncertain{Évidemment} est
alb-équiv.\ à $0$ sur $X$, il est évident qu'il en est de même de
$-\mathfrak{z}$ :} Prouvons que si $\mathfrak{z}, \mathfrak{z}'$ sont
alb-équiv.\ à $0$ sur $X$, il en est de même de $\mathfrak{z} -
\mathfrak{z}'$. Soient $T, T'$ les \uncertain{variétés} de paramètres comme
dans 1.1.\ correspondant à $\mathfrak{z}, \mathfrak{z}'$. Soit $T_{1} = T
\times T'$.

\page{25}
Alors $Z, Z'$ ont des images inverses $Z_{1}, Z'_{1}$ sur $X \times T_{1}$,
et si $\mathfrak{t}, \mathfrak{t}'$ sont les $0$-cycles de $T, T'$ comme dans
1.1., — par les \ill{} \uncertain{on} $\mathfrak{t}_{1} = \mathfrak{t}
\times a'$, $\mathfrak{t}'_{1} = a \times \mathfrak{t}'$ de $T \times T'$.
\note{les points $a$, $a'$ de $T$, $T'$ ne sont pas introduits sur la page}
Notons que
\marginal{Lemme à dégager}
\[
  S_{T \times T'}(\mathfrak{t} \times a') = S_{T}(\mathfrak{t}) = 0, \qquad
  S_{T \times T'}(a \times \mathfrak{t}') = S_{T'}(\mathfrak{t}') = 0
\]
donc $\mathfrak{t}_{1}$ et $\mathfrak{t}'_{1}$ sont [de degré $0$ et]
\struck{alb-équiv.} tels que leurs $S_{T \times T'}$ sont nuls. D'autre part,
on voit de suite que
\[
  Z_{1}(\mathfrak{t}_{1}) = Z(\mathfrak{t}), \qquad Z'(\mathfrak{t}'_{1}) =
  Z'(\mathfrak{t}'),
\]
\struck{donc on est} d'autre part
\[
  Z_{1}(\mathfrak{t}'_{1}) = Z(\mathrm{pr}_{1*}(\mathfrak{t}'_{1})) = 0
\]
car $\mathrm{pr}_{1*}(\mathfrak{t}'_{1}) = 0$, sous réserve que $a$
\uncertain{ait été} choisi tel que $Z(a)$ soit défini. De même, si $a'$ est
choisi convenablement, $Z'_{1}(\mathfrak{t}_{1}) = 0$. On posera donc
\[
  Z_{2} = Z_{1} + Z'_{1}, \qquad \mathfrak{t}_{2} = \mathfrak{t}_{1} +
  \mathfrak{t}'_{1},
\]
et on aura
\[
  Z_{2}(\mathfrak{t}_{2}) = \mathfrak{z} + \mathfrak{z}' .
\]
\note{il annonçait $\mathfrak{z} - \mathfrak{z}'$ ; le calcul donne
$\mathfrak{z} + \mathfrak{z}'$, le passage à $-\mathfrak{z}$ ayant été biffé
plus haut}
\struck{\ill{}} Donc $N(X)$ est bien un s-groupe de $X$.

\note{croquis en marge, à hauteur de la définition de $Z_2$ :}

\begin{tikzcd}
X \times T & X \times T \times T' \arrow[l] \arrow[d, no head] \\
T \arrow[u, no head] & T \times T' \arrow[l]
\end{tikzcd}

\page{26}
Le fait qu'il soit gradué est immédiat. Il contient le groupe des cycles
lin.\ équiv.\ à $0$, comme il résulte aussitôt des définitions, et du fait que
sur $T = \mathbb{P}^{1}$, tout $0$-cycle $\mathfrak{t}$ de degré $0$ satisfait
$S_{T}(\mathfrak{t}) = 0$. \struck{Ainsi $N(X)$ est bien un s-groupe de $X$,
contenant}

Pour prouver (iii), notons $\mathfrak{z}$ un cycle de la classe envisagée
dans $A^{*}(Y)$, et considérons un $Z$ sur $Y \times T$ et un $\mathfrak{t}$
comme dans 1.1., donc $\mathfrak{z} = Z(\mathfrak{t})$. Si nous changeons $Z,
\mathfrak{t}$ dans leur classe d'équivalence linéaire [en $Z', \mathfrak{t}'$],
$Z'(\mathfrak{t}')$ sera linéairement équivalent à $\mathfrak{z}$, et le
$f^{*}$ \uncertain{est le même} pour $\mathfrak{z}$ et $\mathfrak{z}'$. Il
suffit maintenant de prouver qu'il existe $Z', \mathfrak{t}'$ tels que
$(f \times \mathrm{id}_{T})^{*}(Z') = Z'_{1}$ et $Z'(\mathfrak{t}')$,
\struck{$Z'_{1}(\mathfrak{t}')$} soient définis ; en effet, en vertu de la
théorie des intersections, on aura alors
\[
  f^{*}(Z'(\mathfrak{t}')) = Z_{1}(\mathfrak{t}'),
\]
le premier membre \struck{représentant et dans la} représente
$f^{*}(\mathfrak{z})$ dans $A^{*}(X)$, \struck{le deuxième est}

\note{croquis en marge :}

\begin{tikzcd}
X \arrow[r] & Y \\
X \times T \arrow[r] \arrow[d] & Y \times T \arrow[dl] \\
T &
\end{tikzcd}

\page{27}
\marginal{à vérifier}
Nous laissons ici le détail de la vérification de l'existence de $Z',
\mathfrak{t}'$ !

On prouve (ii) par \struck{la} la même méthode que (i), \struck{\ill{}} en se
ramenant au cas où $\mathfrak{z}, \mathfrak{z}'$ sont définis en termes des
mêmes $T, \mathfrak{t}$ et deux cycles $Z, Z'$ sur $T$. On [peut] se ramener
aussi à (iii) et à l'énoncé analogue à (ii) par les produits
\emph{cartésiens}, ce qui dispense de faire au moins bouger les cycles.

La démonstration de (iv) se fait par l'\uncertain{AQT}. Celle de (v) est
triviale sur les définitions (car \ill{} le formalisme de l'équiv.\ alg.)

\struck{Théorème} \emph{Scholie 1.3.} Les $A'^{*}(X)$, pour $X$ variable, ont
les \struck{\ill{}} propriétés fonctorielles \uncertain{familières} que les
$A^{*}(X)$. En particulier, on a la formule de projection. Dans les cas où
équiv.\ linéaire $=$ équiv.\ alg.\ sur $X$ (p.\ ex.\ lorsque
\uncertain{ponctuel}, les $\mathbb{P}^{n}$, les \struck{\ill{}} variétés de
drapeaux) elles sont $=$ équiv.\ alb.
\note{le scholie est marqué d'un trait vertical dans la marge gauche}

\page{28}
\marginal{Ici il semble qu'on utilise le fait que dans 1.1.\ on prend $T$
projectif}
\emph{Théorème 1.4.} Soit $X$ projective lisse connexe de dim $n$ sur $k$.
\struck{Alors pour les cycles de dim $0$ et pour les cycles de codim $1$,
l'équiv.}

(i) Soit $\mathfrak{t}$ un $0$-cycle sur $X$. Alors $\mathfrak{t}$ est
alb-équiv.\ à $0$ ssi $\mathfrak{t}$ est de degré $0$, et
$S_{X}(\mathfrak{t}) = 0$.

(ii) Soit $D$ un diviseur sur $X$. Alors $D$ est alb-équiv.\ à $0$ ssi $D$
est lin.\ équivalent à $0$.

\emph{Remarque 1.5.} On ignore, [déjà] pour les $0$-cycles, si l'équiv.\
d'Alb.\ est identique à l'équivalence linéaire (vrai pour tous les cycles
???).

\emph{Démonstration} (i) Si $\mathfrak{t}$ est de degré $0$ et
$S_{X}(\mathfrak{t}) = 0$, — $\mathfrak{t}$ alb-équiv.\ à $0$ par définition
[prendre $T = X$, $Z = $ diagonale de $X$, donc $\mathfrak{t} =
Z(\mathfrak{t})$]. L'assertion en sens inverse résulte de [ , 1.2.].
\note{la référence est laissée en blanc par lui}
Pour (ii), on sait déjà que lin.\ équiv.\ à $0$ $\Longrightarrow$ alb-équiv.\
à $0$ ; prouvons le sens inverse. On doit prouver que si

\page{29}
$Z$ est un diviseur sur $X \times T$ et $\mathfrak{t}$ sur $T$ est de degré
$0$ et tel que $S_{T}(\mathfrak{t}) = 0$, alors $Z(\mathfrak{t})$ est lin.\
équiv.\ à $0$. Or on a $Z(\mathfrak{t}) = u(S_{T}(\mathfrak{t}))$, où $u :$
\struck{$\mathrm{Alb}$} $A_{T} \to \mathrm{Pic}^{0}_{X}$ est l'hom défini
par $Z$.

\emph{Théorème 1.6.} Soient $\mathfrak{z}, \mathfrak{z}'$ deux classes dans
$A'^{*}(X)$ alg.\ équiv.\ à $0$, $X$ étant \struck{lisse projective}. Alors
$\mathfrak{z}\mathfrak{z}' = 0$. [Si $X$ est projective, alors il suffit que
l'un des deux $\mathfrak{z}$ soit alg.\ équiv.\ à $0$, et l'autre \emph{num.}\
équiv.\ à $0$ ??]
\marginal{Voir s'il ne suffit pas que $\mathfrak{z}$ \ill{} $\mathfrak{z}'$
soit \uncertain{num.\ équiv.} \ill{} \ill{}}
\marginal{des conditions à en venir !}
\note{les deux notes marginales sont écrites en travers, dans la marge
gauche ; la première, sur des lignes tracées, est reliée par une double
flèche au crochet du texte}
\struck{La} [La démonstration] résulte formellement de la définition [via
variétés $T$ \emph{projectives} lisses connexes] de « alg.\ équiv.\ à $0$ »
[et du fait que si $T, T'$ sont deux variétés [telles] \struck{et} (lisses
connexes suffit) et $x, y$ sont deux points de $T$, $x'$ et $y'$ deux pts de
$T'$, alors
\[
  \mathfrak{t} = \bigl((x) - (y)\bigr) \times \bigl((x') - (y')\bigr) =
  (x, x') - (x, y') - (y, x') + (y, y')
\]
est tel que \struck{\ill{}} $S_{T \times T'}(\mathfrak{t}) = 0$.
\marginal{NB. On ignore si cet $\mathfrak{t}$ est lin.\ équiv.\ à $0$, même
si $T, T'$ sont des courbes !}
[Il y a intérêt à transformer le \struck{\ill{}} [\uncertain{lemme}] en un
énoncé sur l'opération $\times$]

\page{30}
\subsection*{2. Familles de classes de cycles mod équiv.\ d'Alb.}

\note{« mod équiv.\ d'Alb. » remplace, dans l'interligne, un « à équivalence »
biffé}

\emph{Théorème 2.1.} Soient $X^{n}$ et $Y^{m}$ comme d'habitude et $Z \in
A^{i}_{\mathrm{alb}}(X \times Y)$, d'où une \struck{loi compo} [application]
$\mathfrak{t} \mapsto Z(\mathfrak{t})$ de \struck{$A_{0}(Y)$}
$A_{0\,\mathrm{alb}}(Y)$ dans $A^{i}_{\mathrm{alb}}(X)$. \struck{Il existe une
application et une seule de} [$A_{0\,\mathrm{alb}}(Y) \simeq
\mathrm{Alb}^{*}_{Y}(k)$, dans $A^{*}_{\mathrm{alb}}(X)$]
\struck{telle que pour tout $\mathfrak{t} \in \mathrm{Alb}_{0}(Y)$} ; or
$A_{0\,\mathrm{alb}}(Y) \overset{S_Y}{\simeq} A^{*}_{Y}(k)$, d'où un hom de
groupes
\[
  u^{*}_{Z} : A^{*}_{Y}(k) \longrightarrow A^{i}_{\mathrm{alb}}(X),
\]
\note{sous $A^{*}_{Y}(k)$, entouré, un « $\|$ Alb » au crayon}
\struck{égal à la restriction à $A^{(n)}_{Y}(k)$} est défini par la condition \struck{$S_Y$}
\[
  Z(\mathfrak{t}) = u_{Z}(S(\mathfrak{t}))
\]
pour tout $\mathfrak{t}$ (resp.\ pour tout $\mathfrak{t}$ de degré $0$),
(\uncertain{i.e.\ en négligeant} la restriction de degré $0$).
\emph{De plus}, si on désigne l'hom \emph{induit} sur $A_{Y}(k) =
A^{0}_{Y}(k)$
\[
  u_{Z} : A_{Y}(k) \longrightarrow A^{i}_{\mathrm{alb}}(X) ,
\]
on a ce qui suit : \struck{il existe une base de}
\begin{enumerate}
  \item[(i)] Il ne dépend que de la classe de $Z$ dans
  $A^{i\,\tau}(X \times Y)$ \marginal{[À \uncertain{réserver} pour plus
  tard]}
  \item[(ii)] Il existe une \struck{base} [isogénie] \struck{isogénie} de VA
  $\varphi : A' \to A_{Y}$ et un cycle $Z' \in A^{i}_{\mathrm{alb}}($
\end{enumerate}
\note{l'indice de $A^{i\,\tau}$ est biffé et illisible. La fin de (ii) est
écrite dans un cadre, en bas à gauche du feuillet, et reprise en haut du
feuillet 31}

\page{31}
de VA
\[
  \varphi : A' \to A_{Y}
\]
et un $Z' \in A^{i}_{\mathrm{alb}}(X \times A')$,
\struck{dont} tel que $Z'((0)_{A'}) = 0$ \marginal{\uncertain{diviseur} de
degré $1$ concentré en $0$} (ce qui revient au même, tel que l'application
$v_{Z'} : a' \mapsto Z'(a')$ de $A'(k)$ dans
\struck{$A^{i}_{Z}$} $A^{i}_{\mathrm{alb}}(X)$ soit un hom de groupes) tel
que
\[
  u_{Z} \circ \varphi(k) = v_{Z'}
\]
\note{dans $A^{i}_{\mathrm{alb}}(X \times A')$, un premier $X$ est biffé}

\emph{Dém.} Supposons qu'il existe un \struck{$Z$} [cycle] $U$ sur $Y \times
A_{Y}$, de codim $m$, tel que pour tout $a \in A_{Y}$, on ait
$S_{Y}(U(a)) = a$. Alors la « famille algébrique de diviseurs »
\uncertain{comprise} \ill{} \ill{}
\[
  Z(U(a)) = Z'(a)
\]
\ill{} \ill{} \ill{} pour tout $\mathfrak{t} \in$ \struck{\ill{}}
$A_{0\,\mathrm{alb}}(Y)$ [de degré $0$,] on aura
\struck{$Z'(a)$}
\[
  Z(\mathfrak{t}) = Z'(S(\mathfrak{t}))
\]
\struck{\ill{} pour $\mathfrak{t}$ \ill{} de \ill{} pour \ill{}}.
Car par le \ill{}, \struck{pour \ill{}} $U(a)$, dont la formule vérifiée
\struck{\ill{} équivalent}
\[
  Z(U(a)) = Z'(S(U(a))) = Z'(a) \qquad \text{O.K.}
\]
Mais à priori on ne sait pas si on peut

\page{32}
trouver $U$, mais on peut trouver $U$ [comme dessus] de façon que l'on ait
\[
  S_{Y}(U(a)) = na
\]
et alors, posant $Z'(a) = Z(U(a))$, la relation
\[
  Z'(S_{Y}(\mathfrak{t})) = n\, Z(\mathfrak{t}) \qquad \text{$\mathfrak{t}$ de
  degré $0$ sur $Y$.}
\]
\struck{(car on peut écrire $n S(\mathfrak{t}) = U(a)$ si $a =
S_{Y}(\mathfrak{t})$,} \struck{$n\mathfrak{t}$ de la forme $U(a')$, car si $a'
= S_{Y}(\mathfrak{t})$,} \struck{on aura $S_{Y}(U(a')) = na' =
S_{Y}(n\mathfrak{t})$ et} \struck{$U(a') = n\mathfrak{t}$, et comme on aura
$a' = na$ pour $U(a)$} \struck{\ill{} \ill{} \ill{}}
\note{passage biffé par un cadre et quatre traits obliques ; lecture
partielle}
Pour le voir, soit $a = S_{Y}(\mathfrak{t})$, on aura
\struck{$U(a) = n\mathfrak{t}$} d'où \struck{\ill{}}
\[
  S_{Y}(U(a)) = na = S_{Y}(n\mathfrak{t}),
\]
\[
  U(a) = n\mathfrak{t}
\]
d'autre part on aura
\[
  Z'(S_{Y}(U(a))) = n\, Z(U(a)), \qquad Z'(S_{Y}(U(a))) = Z'(na), \quad
  n\,Z(U(a)) = n\,Z'(a) .
\]
\note{les deux dernières égalités sont écrites sous la précédente, reliées
par des signes $\|$}
\struck{d'où, comme $U(a) = n\mathfrak{t}$} Or tout élément de
$A_{0\,\mathrm{alb}}(Y)$ de degré $0$ est \struck{$Z'(S_{Y}(\mathfrak{t})) =
n^{2} Z(\mathfrak{t})$} de la forme $n\mathcal{L}$ (car $A_{Y}(k)$ est
divisible !), donc appliquant ce qui précède à $\mathcal{L}$ au lieu de
$\mathfrak{t}$, on obtient
\[
  Z'(S_{Y}(\mathfrak{t})) = n\, Z(\mathfrak{t}) .
\]
On peut comme toujours négliger la multiplication par $n$.

\page{33}
On prouve (ii). Pour (i), on note qu'on sait déjà que si $Z$ est alg.\
équiv.\ à $0$, alors l'hom correspondant $u : A_{Y}(k) \to
A^{i}_{\mathrm{alb}}(X)$ est nul (th.\ 1.6.), donc si $Z$ est
$\tau$-équivalent à $0$, alors l'hom $u$ est tel que $mu = 0$ pour quelque
$m > 0$, donc $u = 0$ car $A_{Y}(k)$ est un divisible.

\emph{Corollaire \uncertain{2.2}.} Pour qu'un cycle algébrique de codim $i$
sur $X$ soit alg.\ équiv.\ à $0$, il f.\ et s.\ qu'il soit
Albanese-équivalent à un diviseur de la forme $Z(a_{0})$, où \struck{$A$ est
une} $Z$ est un cycle de codim $i$ sur une \struck{VA} $A$, \uncertain{plutôt}
$A \times X$, $A$ une variété abélienne, tel que \struck{$\mathrm{Alb}$} $a
\mapsto Z(a)$ : $A \to A^{i}_{\mathrm{alb}}(X)$ soit un hom de groupes.
\note{le numéro du corollaire est surchargé ; suivent 2.3 et 2.4}

\emph{Corollaire 2.3.} Pour tout entier $n > 0$, le s-groupe [de
$A^{i}_{\mathrm{alb}}(X)$, formé des classes] $P^{i}(X)$ des
\struck{classes d'équivalence, à équiv.\ d'Alb.} alg.\ équiv.\ à $0$, est
$n$-\emph{divisible}.

\emph{Corollaire 2.4.} Soient $Z, Z' \in A^{*}_{\mathrm{alb}}(X)$, dont l'un
est alg.\ équiv.\ à $0$ et l'autre $\tau$-équiv.\ à $0$. Alors $ZZ' = 0$.

\page{34}
\emph{Définition 2.5.} Soient $X, Y$ … et considérons une application
\[
  \varphi : Y(k) \longrightarrow A^{i}_{\mathrm{alb}}(X)
\]
On dit que cette application est \emph{algébrique} si
\begin{enumerate}
  \item[1°)] L'hom \struck{$Z$} $\mathrm{Cyc}_{0}(Y) \to
  A^{i}_{\mathrm{alb}}(X)$ correspondant se factorise par
  $A_{0\,\mathrm{alb}}(Y) \simeq A^{*}_{Y}(k)$, i.e.\ \struck{s'annule} est
  nulle sur les $\mathfrak{t}$ de degré $0$ tels que $S_{Y}(\mathfrak{t}) =
  0$
  \item[2°)] L'hom induit $\varphi : A^{0}_{Y}(k) \to
  A^{i}_{\mathrm{alb}}(X)$ est tel que pour $N > 0$ convenable, $N\varphi$
  [algébrique] puisse être induit par un cycle $Z \in
  A^{i}_{\mathrm{alb}}(X \times A_{Y})$
\end{enumerate}
\note{au-dessus de « $N>0$ convenable, $N\varphi$ », dans l'interligne :
$N\varphi(x) = \varphi(Nx)$. En marge gauche, à hauteur de 1°) :}

\begin{tikzcd}
Y \arrow[r] & A^{1}_{Y} \arrow[r, "N"] & A^{N}_{Y}
\end{tikzcd}

\emph{Prop.\ 2.6.} Si $\varphi$ est définie par un $Z \in
A^{i}_{\mathrm{alb}}(X \times Y)$, elle est \emph{algébrique}.

C'est 2.1. Voir la généralisation en 2.7.\ (ii)

\emph{Prop.\ 2.7.}
\begin{enumerate}
  \item[(0)] Forment un s-groupe. Si $N\varphi$ est alg., $\varphi$ l'est.
  \item[(i)] Fonctorialité en $Y$. \struck{Plus généralement}
  \item[(ii)] Fonctorialité en $X$, plus précisément on a un hom
  $A^{i}_{\mathrm{alb}}(X) \to A^{j}_{\mathrm{alb}}(X')$ défini par une
  classe de correspondances algébriques à coeff.\ \struck{rationnels}
  entiers.
\end{enumerate}
\note{l'exposant du but, à la seconde flèche, est peu net}

\page{35}
N.B. Fixons un pt $y_{0}$ de $Y$. Alors le groupe des familles algébriques
$Y(k) \to A^{i}_{\mathrm{alb}}(X)$ est somme de
\[
  A^{i}_{\mathrm{alb}}(X) \times \mathrm{Homalg}(A_{Y},
  A^{i}_{\mathrm{alb}}(X)) ,
\]
où le deuxième facteur désigne les familles algébriques qui sont des hom de
groupes, i.e.\ nulles à l'origine. On peut donc se limiter à étudier les hom
« alg. » $A(k) \to A^{i}_{\mathrm{alb}}(X)$ [où $A$ est une VA], et
interprétant $A$ comme
\[
  A_{0\,\mathrm{alb}}(A) = A^{\alpha}_{\mathrm{alb}}(A)^{(\text{alg.\ équiv.\
  }0)} \qquad (\alpha = \dim A)
\]
on peut considérer ce concept comme un cas particulier des \emph{hom
algébriques}
\marginal{$A(k) \simeq P^{\alpha}(A) \to P^{i}(X)$}
\[
  \varphi : P^{j}(X') \longrightarrow P^{i}(X) .
\]
Un hom de groupes $\varphi$ est dit « algébrique » s'il existe un entier $N >
0$ tel que $N\varphi$ soit définissable par un cycle algébrique sur $X \times
X'$. (Alors $\varphi$ sera définissable par un cycle alg.\ à coeff.\
\emph{rationnels} — mais attention, un cycle à coeff.\ rationnels
\emph{quelconque} ne définit pas en général un hom $P^{j}(X') \to P^{i}(X)$,
à moins que les degrés soient les bons …). Ces hom se \emph{composent}

\page{36}
\marginal{?}
Une question plus subtile est s'ils peuvent s'inverser, quand ce sont des
isomorphismes …
\note{ces deux lignes sont marquées d'un trait vertical dans la marge}

\emph{Dualité} entre $P^{i}(X)$ et \struck{$P^{n-i}$} $P^{j}(X)$ [pour $i + j
= n + 1$] : Si $A, B$ \struck{sont deux} [des] VA, \struck{dont la dualité est}
accouplées par une dualité parfaite ($A \simeq B^{*}$, $B \simeq A^{*}$),
alors il y a correspondance biunivoque entre
$\mathrm{Homalg}(A, P^{i}(X))$ et $\mathrm{Homalg}(P^{j}(X), B)$ :

En effet, si on a un hom $A \to P^{i}(X)$, défini par un $Z \in
A^{i}_{\mathrm{alb}}(X \times A)$, alors ${}^{t}Z$ définit un hom $P^{j}(X)
\to P^{i+j-n}(A) = P^{1}(A) \simeq B$, [i.e.] $P^{j}(X) \to B$.
\struck{Notation} Cependant, il \uncertain{s'ensuivra} si on \ill{}
\struck{tq \ill{}. Plus généralement} il faut vérifier que $P^{j}(X) \to B$ ne
dépend que de l'\emph{application} $A \to P^{i}(X)$, pas du cycle $Z$.
\note{ces deux lignes sont marquées d'un triple trait vertical dans la
marge}

Question de transposition de façon plus générale, pour $P^{i}(X) \to
P^{j'}(X')$ en $P^{i'}(X') \to P^{j}(X)$. Cf.\ plus bas …

\page{37}
\subsection*{3. Homomorphismes $\ell$-adiques associés à une famille abélienne de cycles}

Soit $A$ une V.A., $Z \in A^{i}_{\mathrm{alb}}(X \times A)$, d'où
\[
  u_{Z} : A(k) \longrightarrow A^{i}_{\mathrm{alb}}(X)
\]
mais aussi
\[
  a({}^{t}Z) : H^{2(n-i)+1}(X)(n-i) \longrightarrow H^{1}(A)
\]
(coefficients $\ell$-adiques \struck{\ill{}} $\mathbb{Z}_{\ell}$). Je dis que
mod torsion (i.e.\ à coeff.\ $\mathbb{Q}_{\ell}$ !) cet homomorphisme ne
dépend que de $u_{Z}$, pas de $Z$, \struck{donc} et de façon précise que :

\marginal{? A-t-on une réciproque}
\emph{Théorème 3.1.}
\[
  u_{Z} = 0 \Longrightarrow a({}^{t}Z) : H^{2(n-i)+1} \to H^{1} \text{ est
  nul}
\]
\marginal{dém.\ pas faite}
\note{au-dessus de l'implication, un « ? » ; à sa gauche, un chevron « < »
qui semble poser la réciproque. « dém.\ pas faite » est entouré au crayon et
relié au mot « Théorème »}

\emph{Démonstration.} Choisissons une « courbe ample » $C \to A$
\ill{} $C$ sur $A$, \struck{disons} d'où $Z' \in
A^{i}_{\mathrm{alb}}(C \times X)$ \struck{$C \times X$}
[$X \times C \to X \times A$] \struck{$A_{C} \to A$} par conséquent
\[
  u_{Z'} : C(k) \longrightarrow A^{i}_{\mathrm{alb}}(X)
\]
\[
  a({}^{t}Z') : H^{2(n-i)+1}(n-i) \longrightarrow H^{1}(C)
\]
\note{un « $(X)$ » est biffé après $H^{2(n-i)+1}$}
et comme $H^{1}(A) \to H^{1}(C)$ est injectif, $C(k) \to \mathrm{Alb}$,
surjectif, — .

\page{38}
\[
  u_{Z} = 0 \Longleftrightarrow u_{Z'} = 0, \qquad a({}^{t}Z) = 0
  \Longleftrightarrow a({}^{t}Z') = 0 .
\]
On est donc ramené à la question analogue, avec $A$ remplacé par $C$. Or
\struck{tout} \uncertain{supposons} que la décomposition
\[
  H^{N}(X \times C) \simeq H^{N}(X) \otimes H^{0}(C) + H^{N-1}(X) \otimes
  H^{1}(C) + H^{N-2}(X) \otimes H^{2}(C)
\]
est algébrique, comme \ill{} on \uncertain{vérifie} aisément. Appliquant ceci
avec $N = 2i$, et notant que un cycle algébrique « appartenant » à l'un des
deux facteurs extrêmes donne un $u_{Z'}$ \emph{et} un $a({}^{t}Z)$ nul, on est
ramené au cas d'un cycle « appartenant » au facteur médian, i.e.\ dont la
\emph{projection} sur $X$, et la « restriction » à $X$, est nulle.
\note{le feuillet s'arrête aux deux tiers ; la démonstration n'est pas
achevée ici}

\section{Astuce de Tate sur les cycles décomposables}

\note{titre de sa main, seul sur le feuillet 39, laissé sans \emph{page}
ici comme une couverture. Le feuillet 40, qui suit, est d'une encre plus
noire, sur un papier jauni}

\page{40}
Soit $A$ une VA sur $k$ alg.\ clos, d'où module $\mathbb{T}_{\ell}(A)$
\struck{\ill{}} sur la pseudo-alg.\ de Lie \struck{\ill{}}
$\mathfrak{g}(k) = \mathfrak{g}$. Soit $K$ une clôture alg.\ de
$\mathbb{Q}_{\ell}$, posant
\[
  \mathbb{T}_{K}(A) = \mathbb{T}_{\ell}(A)
\]
\note{sic : sans $\otimes K$}
de sorte que $\mathbb{T}_{K}(A)$ est un module sur $\mathfrak{g}
\otimes_{\mathbb{Q}_{\ell}} K = \mathfrak{g}_{K}$. Supposons

a) La \struck{\ill{}} conj.\ de Tate vérifiée pour $H^{2}(A,
\mathbb{Q}_{\ell}(1))$
\[
  H^{2}(A, \mathbb{Q}_{\ell}(1))^{\mathrm{alg}} = H^{2}(A,
  \mathbb{Q}_{\ell}(1))^{\mathfrak{g}}
\]

b) $\mathfrak{g}_{K}$ opérant sur \struck{$\mathfrak{g}$} $\mathbb{T}_{K}(A)$
est « pseudo-algébrique » [i.e.\ les valeurs propres des Frobenius
\struck{sont} de \ill{} $q$ sont $\varepsilon q^{-1/2}$, $\varepsilon$ une
racine de l'unité …] \emph{ou} $\mathbb{T}_{K}(A)$ se décompose en deux
morceaux $V_{1} \times V_{2}$, tels que chacun soit propre pour « presque
tous » les Frobenius, avec des valeurs propres différentes pour presque tous
les Frobenius.
\marginal{Par \uncertain{Tate}, la première \ill{} est un cas particulier de
la seconde}

[Ces] conditions sont vérifiées si $A$ est une [puissance]
\struck{produit} d'une courbe elliptique (définissable) sur un corps fini
$\mathbb{F}_{q}$, comme on voit en notant que a) [résulte] \struck{\ill{} se
ramène au cas d'une puissance cartésienne et que pour $C$, a) est} a été
vérifié par Mumford (comme on sait
\note{la phrase se poursuit au-delà du feuillet}

\end{document}
